% ============================================================================
\chapter{Training the Logic Gate Network on Clause Outputs}
\label{ch:lgn-training}
% ============================================================================

This chapter describes how to train a Logic Gate Network (LGN)~\cite{petersen2022difflogic} using
the clause outputs from a pre-trained Tsetlin Machine~\cite{granmo2018tsetlin} as input. This
is Phase~2 of the hybrid architecture: the Tsetlin Machine is frozen,
and the LGN learns to combine its clause outputs into a better
classifier.


% ----------------------------------------------------------------------------
\section{The Training Setup}
\label{sec:lgn-setup}
% ----------------------------------------------------------------------------

After Phase~1 (Tsetlin Machine training), we have:
\begin{itemize}
  \item A trained Tsetlin Machine with $m$ clauses per class and
    $K = 10$ classes, giving $m \times K$ total clauses.
  \item For each training image, a binary vector of clause outputs:
    $\bvec{c} \in \{0, 1\}^{m \times K}$.
\end{itemize}

The LGN takes $\bvec{c}$ as input and produces a classification
decision. The goal is to learn a Boolean circuit that achieves higher
accuracy than the Tsetlin Machine's~\cite{granmo2018tsetlin} simple vote-and-argmax.

\begin{figure}[htbp]
\centering
\begin{tikzpicture}[
  box/.style={draw, rounded corners, minimum width=2.4cm,
    minimum height=1cm, font=\small\sffamily, align=center,
    line width=0.6pt},
  arr/.style={-{Stealth[length=4pt]}, thick},
]
  \node[box, fill=bitgreen!10] (tm) at (0, 0)
    {Tsetlin Machine\\(frozen)};
  \node[box, fill=bitorange!10] (lgn) at (5, 0)
    {Logic Gate Network\\(training)};
  \node[box, fill=bitpurple!10] (out) at (10, 0)
    {Classification\\output};

  \draw[arr] (-3, 0) -- (tm)
    node[midway, above, font=\scriptsize\sffamily] {image};
  \draw[arr] (tm) -- (lgn)
    node[midway, above, font=\scriptsize\sffamily] {clause bits};
  \draw[arr] (lgn) -- (out)
    node[midway, above, font=\scriptsize\sffamily] {gate outputs};
\end{tikzpicture}
\caption{Phase~2 training: the Tsetlin Machine is frozen and provides
  binary clause outputs. The LGN trains on these outputs using
  gradient descent with the softmax relaxation.}
\label{fig:lgn-setup}
\end{figure}


% ----------------------------------------------------------------------------
\section{Precomputing Clause Outputs}
\label{sec:precompute}
% ----------------------------------------------------------------------------

Since the Tsetlin Machine is frozen during Phase~2, we can precompute
all clause outputs for the entire training set. This is a one-time
cost that dramatically speeds up LGN training.

\begin{engineernote}
For 60,000 training images with 5,000 clauses (500 per class $\times$
10 classes), the precomputed dataset is:
\[
  60{,}000 \times 5{,}000 \text{ bits}
  = 300{,}000{,}000 \text{ bits}
  = 37.5 \text{ MB}
\]
This fits easily in memory. Each training epoch of the LGN then
reads from this cached binary matrix instead of re-evaluating the
Tsetlin Machine.
\end{engineernote}


% ----------------------------------------------------------------------------
\section{Network Architecture for MNIST}
\label{sec:lgn-arch}
% ----------------------------------------------------------------------------

A concrete LGN architecture for MNIST:

\begin{description}[leftmargin=1em]
  \item[Input layer:] 5,000 bits (clause outputs from the TM).

  \item[Hidden layers:] 4 layers, each with 4,000 nodes. Each node
    has two randomly selected inputs from the previous layer. Each
    node maintains 16 logits for gate selection.

  \item[Output layer:] 10 groups of nodes, one per digit class. Each
    group has some number of output gates. The popcount of each
    group's outputs gives the confidence score for that class.

  \item[Classification:] $\hat{y} = \arg\max_k \; \text{popcount}
    (\text{group}_k)$.
\end{description}

During training, the continuous relaxation produces real-valued
outputs in $[0, 1]$, and we use cross-entropy loss:
\begin{equation}
  \mathcal{L} = -\log \frac
    {\exp(s_y)}
    {\sum_{k=0}^{9} \exp(s_k)}
  \label{eq:lgn-loss}
\end{equation}
where $s_k$ is the (continuous) score for class $k$ and $y$ is the
true label.


% ----------------------------------------------------------------------------
\section{Training Procedure}
\label{sec:lgn-procedure}
% ----------------------------------------------------------------------------

\begin{algorithm}[htbp]
\caption{Training the Logic Gate Network (Phase~2).}
\label{alg:lgn-train}
\begin{algorithmic}[1]
  \State \textbf{Input:} Precomputed clause outputs $\{(\bvec{c}_i, y_i)\}_{i=1}^{60000}$
  \State Initialise logits $\bvec{z}$ randomly (small values near 0)
  \State Set temperature $\tau \leftarrow \tau_{\text{start}}$ (e.g., 5.0)
  \For{epoch $= 1, 2, \ldots, E$}
    \For{each mini-batch $B$ of clause output vectors}
      \State \textbf{Forward pass:}
      \For{each node $i$ in topological order}
        \State $p_i^{(g)} \leftarrow \text{softmax}(\bvec{z}_i / \tau)$
          for $g = 0, \ldots, 15$
        \State $\hat{y}_i \leftarrow \sum_g p_i^{(g)} \cdot
          \TruthTable(g)[a_i, b_i]$
      \EndFor
      \State Compute loss $\mathcal{L}$ using \cref{eq:lgn-loss}
      \Statex
      \State \textbf{Backward pass:}
      \State Compute $\frac{\partial \mathcal{L}}{\partial \bvec{z}}$
        via backpropagation
      \State Update logits: $\bvec{z} \leftarrow \bvec{z} -
        \eta \frac{\partial \mathcal{L}}{\partial \bvec{z}}$
    \EndFor
    \State Anneal temperature: $\tau \leftarrow \tau \cdot \gamma$
      \Comment{$\gamma < 1$, e.g., 0.95}
  \EndFor
  \Statex
  \State \textbf{Discretise:}
  \State For each node $i$: $g_i^* \leftarrow \arg\max_g \; z_i^{(g)}$
  \State Discard logits; fix each node to gate type $g_i^*$
\end{algorithmic}
\end{algorithm}


% ----------------------------------------------------------------------------
\section{What the LGN Learns}
\label{sec:lgn-learns}
% ----------------------------------------------------------------------------

The LGN learns nonlinear combinations of clause outputs that the
Tsetlin Machine's linear voting cannot express. Examples:

\paragraph{Clause conjunction.}
``If clause~17 (top-bar pattern) fires AND clause~42 (bottom-curve
pattern) fires, this is likely a `5'.'' The LGN can implement this
with a single AND gate connecting these two clause outputs.

\paragraph{Clause inhibition.}
``If clause~23 (vertical-stroke pattern) fires AND clause~89
(horizontal-cross pattern) does NOT fire, this is a `1' not a `4'.''
The LGN implements this with a gate of type 2 ($A \wedge \bar{B}$).

\paragraph{Feature hierarchy.}
Through multiple layers, the LGN builds a hierarchy: layer~1 combines
pairs of clause outputs, layer~2 combines those combinations, and so
on. By layer~4, each gate's output depends on up to $2^4 = 16$
original clause outputs, enabling complex decision boundaries.

\begin{keyinsight}
The LGN~\cite{petersen2022difflogic} adds \emph{depth} to the Tsetlin Machine's flat structure.
While the TM evaluates clauses independently and sums the results,
the LGN creates a deep Boolean circuit that can express arbitrary
logical relationships between clause outputs. This compositional
structure is what allows the hybrid to exceed the TM's accuracy.
\end{keyinsight}


% ----------------------------------------------------------------------------
\section{Accuracy Improvement}
\label{sec:lgn-accuracy}
% ----------------------------------------------------------------------------

Typical accuracy progression for the hybrid architecture on MNIST~\cite{lecun1998mnist}:

\begin{table}[htbp]
\centering
\caption{Accuracy at each stage of the hybrid architecture.}
\label{tab:stage-accuracy}
\begin{tabular}{@{}lcc@{}}
  \toprule
  \textbf{Stage} & \textbf{Test accuracy} & \textbf{Improvement} \\
  \midrule
  TM alone (vote + argmax) & 96.5\% & --- \\
  TM + LGN (after Phase 2) & 97.5--98.0\% & +1.0--1.5\% \\
  TM + LGN + LUT (after Phase 3) & 97.3--97.8\% & $-$0.2\% (discretisation) \\
  \bottomrule
\end{tabular}
\end{table}

The LGN consistently adds 1--1.5\% accuracy over the TM's linear
voting. The LUT compilation stage introduces a small accuracy loss
(0.1--0.3\%) from quantising continuous scoring functions, which is
addressed in Part~IV.


% ----------------------------------------------------------------------------
\section{Training Considerations}
\label{sec:lgn-considerations}
% ----------------------------------------------------------------------------

\paragraph{Learning rate.}
The logits are the only trainable parameters. A learning rate of
$\eta = 0.01$ with the Adam optimiser works well.

\paragraph{Number of epochs.}
30--100 epochs are typically sufficient. The annealing schedule should
be calibrated so that $\tau$ reaches near zero by the final epoch.

\paragraph{Random connectivity.}
The random wiring of inputs to nodes is fixed before training and
does not change. Different random seeds produce slightly different
final accuracies (variance of about 0.2--0.3\%). The LGN is
surprisingly robust to the specific connectivity pattern.

\paragraph{Gate type distribution.}
After training, the distribution of gate types is informative:
\begin{itemize}
  \item AND and OR gates dominate (together 40--50\% of nodes).
  \item XOR and XNOR are rare but important (5--10\%).
  \item Constant gates (FALSE, TRUE) appear at 10--15\%, effectively
    pruning unused connections.
  \item Projection gates ($A$ or $B$, ignoring the other input)
    appear at 10--15\%, indicating that one of the node's inputs
    is irrelevant.
\end{itemize}

\begin{engineernote}
The presence of projection and constant gates suggests that the
network is over-parameterised. Post-training pruning (removing
constant gates and collapsing projections) can reduce the gate
count by 20--30\% with no accuracy loss.
\end{engineernote}


% ----------------------------------------------------------------------------
\section{The Discrete Inference Circuit}
\label{sec:inference-circuit}
% ----------------------------------------------------------------------------

After discretisation, the LGN is a fixed Boolean circuit. Evaluating
it on a new input requires:

\begin{enumerate}
  \item Compute clause outputs from the Tsetlin Machine (bitwise AND
    operations).
  \item For each gate in topological order, look up the output from
    the gate's truth table using the two inputs.
  \item Aggregate the final layer's outputs to produce a
    classification.
\end{enumerate}

Each gate evaluation is a single operation: form a 2-bit index from
the inputs, look up the result in a 4-entry table (which fits in a
single byte). With 16,000 gates across 4 layers, the total is 16,000
table lookups---trivial for any processor.

\begin{keyinsight}
The discretised LGN has \emph{zero} floating-point operations. Every
gate is a 1-byte table lookup. Combined with the Tsetlin Machine's
bitwise AND operations, the entire pipeline from pixels to
classification uses only Boolean operations, integer addition, and
small table lookups. This is the all-Boolean inference promise of
the hybrid architecture.
\end{keyinsight}


% ----------------------------------------------------------------------------
\section{Summary}
\label{sec:ch9-summary}
% ----------------------------------------------------------------------------

The Logic Gate Network training (Phase~2) transforms the Tsetlin
Machine's linear classifier into a deep Boolean circuit:

\begin{enumerate}
  \item Freeze the Tsetlin Machine and precompute clause outputs.
  \item Train a multi-layer network of 2-input gates using softmax
    relaxation and temperature annealing.
  \item Discretise: select the most probable gate type for each node.
  \item The result is a pure Boolean circuit that improves accuracy by
    1--1.5\% over the TM alone.
\end{enumerate}

Part~IV addresses the remaining piece: how to handle any scoring
functions (like the final vote aggregation) that still involve
arithmetic beyond Boolean operations, by compiling them into lookup
tables.
